\section{Reference averaging, charges, and field products}\label{sec:reference}

The changes of scale below use conditional expectations in the critical
nearest-neighbor model.  This section records their precise scope and derives
three consequences: an invariant coefficient for each of the energy and spin
directions, the large-distance coefficients of the products of these fields,
and a finite set of boundary tests that controls an averaged sup norm.  These
are deterministic statements.  Estimates for the random lists constructed
from them begin in Section~\ref{sec:norms}.

\subsection{The reference measure and its accepted inputs}

Write $\langle\cdot\rangle_0$ for critical nearest-neighbor plane expectation
on $\Z^2$, at coupling $K_c=\frac12\log(1+\sqrt2)$.  For a max-norm square
$B_R(z)$, let $P_{R,z}$ be conditional expectation over its interior, retaining
the spins of its cutting contour.  Thus $P_{R,z}F$ is a function of cutting
spins whenever $F$ is supported inside the square.  All sup norms in this
section are over every assignment of the spins on the support of the
function; for $P_{R,z}F$ these are the cutting spins.  We abbreviate
$P_{R,0}$ to $P_R$.  Nested cuts compose by the tower property.  In particular,
if a square lies inside a larger reference domain, replacing an observable
inside the square by its conditional average preserves its expectation in
that domain, as well as in the plane.

Put $F^\circ=F-\langle F\rangle_0$, and use this representative for even
functions modulo constants.  An odd function already has plane mean zero.
Our scalar energy and spin dimension are
\begin{equation}\label{ref:energy}
 e(z)=\frac14\sum_{v:\,|v-z|=1}
       \bigl(\sigma_z\sigma_v-\langle\sigma_z\sigma_v\rangle_0\bigr),
 \qquad d=\frac18.
\end{equation}
We use Euclidean distances for field asymptotics and max-norm squares for
conditional averaging.  This distinction changes no buffer condition.

The supplied conformal-universality theorem applies to a fixed deterministic
potential on finite even sets, of fixed finite range, invariant under lattice
translations, reflections, and right-angle rotations.  For a sufficiently
small fixed coefficient multiplying that potential, it supplies a critical
branch, nonzero field factors, and separated bulk spin and energy limits in
smooth simply connected domains with free, plus, or minus boundary, and in
the plane obtained from periodic-square thermodynamic limits
\cite[Theorem~1.1]{ReferenceIsing}.  We use its conditional
averaging, thermal, and marked-source interfaces only for the unperturbed
critical reference model in bulk charts with their prescribed buffers.
In particular, its conditional kernels throughout this paper are $P$;
they are never conditional kernels of a random effective interaction.

Here are the specific operator and calibration inputs.  If $F$ is supported
in $B_w(0)$, set
\begin{equation}\label{ref:disk-response}
 d_R(F)=\langle F\rangle^+_{\mathrm{disk}(4R)}-\langle F\rangle_0.
\end{equation}
For each sufficiently large fixed ratio $R/w$, and all sufficiently large
$w$, the reference estimate gives
\begin{align}
 |d_R(F)|&\le C(w/R)^b\norm{F}_\infty,
 \label{ref:response-input}\\
 \norm{P_RF^\circ}_\infty
   &\le C|d_R(F)|+C(w/R)^p\norm{F}_\infty,
 \label{ref:operator-input}
\end{align}
where
\begin{equation}\label{ref:exponents}
 (b,p)=
 \begin{cases}
  (1,3),&F\text{ is even and invariant under right-angle rotations},\\
  (1,2),&F\text{ is even},\\
  (d,1+d),&F\text{ is odd}.
 \end{cases}
\end{equation}
The constant is uniform over sufficiently large fixed ratios, whereas the
fine-mesh threshold can depend on the ratio.  Ratios converging to a fixed
admissible value are allowed.  The norm bound is uniform over all cutting
configurations \cite[Proposition~4.1]{ReferenceIsing}.

Write $\mathcal O_1(z)=e(z)$ and $\mathcal O_d(z)=\sigma_z$,
and omit $z$ when $z=0$.  There are nonzero
constants $\kappa_b$ such that
\begin{equation}\label{ref:calibration}
 R^b d_R(\mathcal O_b)\longrightarrow\kappa_b.
\end{equation}
Moreover the averaged reference fields satisfy
\begin{equation}\label{ref:line-norms}
 c\le\norm{w^bP_w\mathcal O_b}_\infty\le C,
 \qquad
 c(w/R)^b\le |d_R(w^bP_w\mathcal O_b)|\le C(w/R)^b.
\end{equation}
Here and below energy fields are plane-centered.  These are the calibration
and field-line estimates of
\cite[Lemmas~4.5 and~5.3 and Equation~(5.18)]{ReferenceIsing}.  The exact identity
behind the last reference is
\begin{equation}\label{ref:exact-line}
 B^bP_{Bw}\bigl(w^bP_w\mathcal O_b\bigr)
          =(Bw)^bP_{Bw}\mathcal O_b.
\end{equation}
No convergence rate for primary fields is used in these inputs.

\subsection{Invariant energy and spin coefficients}

The disk response in \eqref{ref:operator-input} depends on the observation
scale.  To use a single thermal coordinate through successive cuts, we first
replace that response by a linear functional invariant under all cuts.

\begin{lemma}[Invariant charges]\label{ref:charges}
There are linear functionals $\ell$ on bounded local even functions modulo
constants and $\ell_s$ on bounded local odd functions with
\[
 \ell(e(z))=1,\qquad \ell_s(\sigma_z)=1.
\]
They commute with translations and square symmetries and are invariant under
reference conditional averaging whenever the cut encloses the support.
For $b=1$ write $\ell_b=\ell$, and for $b=d$ write $\ell_b=\ell_s$; the subscript in
$\ell_s$ means spin.  If $F$ has the corresponding parity and support in a
square of radius $w$ above a fixed threshold, then
\begin{equation}\label{ref:charge-bound}
 |\ell_b(F)|\le Cw^b\norm{F}_\infty,
 \qquad
 \ell_b(F)=\kappa_b^{-1}\lim_{R\to\infty}R^b d_R(F).
\end{equation}
The disk in this limit may be centered at any fixed lattice point.  In the
even case $\ell_b$ annihilates constants.
\end{lemma}

\begin{proposition}[Averaging and charge cancellation]\label{ref:operator}
Fix $\nu>0$ smaller than $p-b$ in the applicable row of
\eqref{ref:exponents}.  For all sufficiently large $w$ and $R\ge Cw$, a
function $F$ of that parity, supported in $B_w(z)$, satisfies
\begin{align}
 \norm{P_{R,z}F^\circ}_\infty
   &\le C(w/R)^b\norm{F}_\infty,\label{ref:ordinary-decay}\\
 \norm{P_{R,z}F^\circ-\ell_b(F)P_{R,z}\mathcal O_b(z)}_\infty
   &\le C(w/R)^{p-\nu}\norm{F}_\infty.
   \label{ref:improved-decay}
\end{align}
In particular, the second bound is the improved decay bound for a function
of charge zero.  We shall use $\nu=1/4$.  The exponent $3-\nu$ requires
the individual even function to be invariant under right-angle rotations;
invariance of its probability law alone does not suffice.
For a site at fixed positive relative distance from the cutting contour,
\begin{equation}\label{ref:microscopic-decay}
 \norm{P_{R,z}e(y)}_\infty\le CR^{-1},\qquad
 \norm{P_{R,z}\sigma_y}_\infty\le CR^{-d}.
\end{equation}
The first bound also holds for a single plane-centered bond product.
\end{proposition}

\begin{proof}[Proof of Lemma~\ref{ref:charges} and
Proposition~\ref{ref:operator}]
Fix a sufficiently large integer $B$ and put $w_k=B^kw$.  In the appropriate
parity space, with plane centering in the even case, define
\[
 v_k=w_k^bP_{w_k}\mathcal O_b,
 \qquad S_k=B^bP_{w_{k+1}}.
\]
The lines $\R v_k$ are uniformly nonzero and bounded by
\eqref{ref:line-norms}, and $S_kv_k=v_{k+1}$ by
\eqref{ref:exact-line}.  If $F$ is supported in $B_{w_k}$, subtract
$\alpha v_k$ to cancel its response in the disk of radius $4w_{k+1}$.
Equations~\eqref{ref:response-input} and \eqref{ref:line-norms} give
$|\alpha|\le C\norm{F}_\infty$.  Applying
\eqref{ref:operator-input} to this difference proves the quotient bound
\begin{equation}\label{ref:quotient}
 \norm{S_k:\,X_k/\R v_k\longrightarrow X_{k+1}/\R v_{k+1}}
       \le CB^{b-p}.
\end{equation}
Here $X_k$ is the space of bounded functions of the required parity on
$B_{w_k}$, with plane mean zero in the even case and with right-angle
rotation invariance imposed when $p=3$; its norm is the sup norm.
The same estimate holds after subtracting any multiple of $v_k$, so it is
indeed an estimate for the quotient norm.  Combining \eqref{ref:response-input} and \eqref{ref:operator-input}
at the fixed ratio $B$ also bounds $S_k$ uniformly.

Choose uniformly bounded projections onto $\R v_k$, using for example
the response in a disk at one fixed large multiple of $w_k$ and dividing
by its nonzero response on $v_k$.  In the resulting bounded splittings,
\[
 S_k=
 \begin{pmatrix}1&\beta_k\\0&D_k\end{pmatrix},\qquad
 \norm{\beta_k}\le C,\quad \norm{D_k}\le CB^{b-p}=:q.
\]
Enlarge $B$ until $q\le B^{b-p+\nu}<1$.  Writing
$S_{k-1}\cdots S_0F=a_kv_k+u_k$ gives
\[
 \norm{u_k}\le Cq^k\norm{F}_\infty,\qquad
 a_{k+1}-a_k=\beta_ku_k.
\]
Consequently $a_k$ converges to a linear functional $a_\infty(F)$ and
\begin{equation}\label{ref:line-limit}
 |a_\infty(F)|\le C\norm{F}_\infty,
 \quad
 \norm{S_{k-1}\cdots S_0F-a_\infty(F)v_k}_\infty
       \le Cq^k\norm{F}_\infty.
\end{equation}
Define $\ell_b(F)=w^ba_\infty(F)$.  Since the line is transported exactly,
its coefficient is one on $\mathcal O_b$.  Constants have already been
removed.  Dividing \eqref{ref:line-limit} by $B^{kb}$ gives
\eqref{ref:improved-decay} at the radii $w_k$ and gives
\eqref{ref:charge-bound}'s bound.  The construction permits an arbitrary
initial function on the first square; it need not initially be a function
of the cutting contour.

To identify the functional independently of the chosen initial scale,
first divide \eqref{ref:line-limit} by $B^{kb}$ and test the resulting
unscaled decomposition in a plus disk at a fixed sufficiently large
multiple of $w_k$.  Conditional averaging preserves that expectation.
After multiplication by the disk scale to the power $b$, the error is
$O(B^{k(b-p+\nu)})$ and tends to zero because $p-\nu>b$; the line term tends to
$\kappa_b \ell_b(F)$ by \eqref{ref:calibration}.  This proves the limit in
\eqref{ref:charge-bound}.  It also proves that different initial scales and
bounded splittings give the same functional.  For an intermediate radius $R$, stop one scale earlier when necessary,
choosing $k$ so that $B\le R/w_k<B^2$.  The final averaging ratio then
lies in a fixed admissible compact interval.
The input estimates are uniform on such an interval: otherwise a sequence
of failing ratios would have a convergent subsequence, contrary to their
stated converging-ratio version.  One additional cut therefore proves both
decay bounds and the disk limit for every sufficiently large radius.

For a cut enclosing $F$, its disk expectation agrees with that of $F$ in
every disk containing the cut, and its plane expectation also agrees.
The disk-limit formula now proves $P$-invariance.  A fixed shift of the disk
center gives the same limit: apply the decay decomposition about the old
center, whose remainder still has sup norm $O(R^{-p+\nu})$ in a square
inside the shifted disk.  The line term has the same calibrated one-point
limit, by bulk one-point convergence with the insertion displaced by
$O(R^{-1})$ in rescaled coordinates.  This proves translation invariance;
the disk symmetry proves rotation and reflection invariance.

The bounds in \eqref{ref:microscopic-decay} at the center are
\eqref{ref:line-norms}.  For a displaced site, first cut at a square centered
there of radius comparable to $R$ and contained inside the requested cut,
and then use the tower property and contraction in sup norm.  For a single
centered bond, its one-point calibration has the same nonzero scalar limit
as the four-bond average, by the separated reference field theorem for both
bond orientations.  The recurrence proving \eqref{ref:line-norms} uses
\eqref{ref:operator-input} with $p=2$ in this case: at a fixed large ratio,
\[
 (Bw)\norm{P_{Bw}F}_\infty
       \le C+CB^{-1}w\norm{P_wF}_\infty.
\]
Iteration bounds the left side uniformly.  Finally all four incident bonds
have equal charge by square symmetry, and their average has charge one;
thus every centered bond has charge one.
\end{proof}

\subsection{Extracting product coefficients from separated limits}

Normalize the continuum energy $\varepsilon$ so that its plane two-point
function is $|z-w|^{-2}$.  Let $a>0$ be the lattice energy field factor in
this normalization; a simultaneous sign change of $\varepsilon$ changes
the signed energy--spin coefficient but not any expression used below.
Write $a_\sigma>0$ for the spin field factor.  The separated primary-field
limits, including mixed correlations in a plus disk, are the accepted
reference input and \cite[Theorem~1.3]{CHI2023}.  The continuum product
expansions in a fixed domain are
\begin{equation}\label{ref:continuum-ope}
 \varepsilon(z)\varepsilon(0)=|z|^{-2}\one+O(1),
 \qquad
 \varepsilon(z)\sigma(0)
       =\frac{1}{2|z|}\sigma(0)+O(1).
\end{equation}
They hold inside correlation functions, uniformly away from other
insertions and the boundary
\cite[Theorems~6.1--6.2 and Remark~6.4]{CHI2023}.

\begin{lemma}[Energy and energy--spin products]\label{ref:ope}
Uniformly in the direction of $z\in\Z^2$ as $|z|\to\infty$,
\begin{align}
 \langle e(z)e(0)\rangle_0
       &=\frac{a^2+o(1)}{|z|^2},\label{ref:ee-correlation}\\
 \ell\bigl(e(z)e(0)\bigr)&=o(|z|^{-1}),\label{ref:ee-charge}\\
 \ell_s\bigl(e(z)\sigma_0\bigr)^2
       &=\frac{a^2+o(1)}{4|z|^2}.\label{ref:es-charge}
\end{align}
The even charge in \eqref{ref:ee-charge} is evaluated modulo the plane mean.
In addition to Proposition~\ref{ref:operator}, the proof uses only separated
lattice limits and the continuum expansions \eqref{ref:continuum-ope}.
\end{lemma}

\begin{proof}
Set $n=|z|$ and $\zeta=z/n$.  The plane separated energy limit gives
\eqref{ref:ee-correlation}.  To extract the other two coefficients without
a lattice collision estimate, fix a small $\rho>0$ and cut separately in
the disjoint squares of radius $\rho n$ about $0$ and $z$.  Put
\[
 E_0=P_{\rho n,0}e(0),\quad E_z=P_{\rho n,z}e(z),\quad
 S_0=P_{\rho n,0}\sigma_0,
\]
and
\[
 \Phi_n=n^2(E_zE_0-\langle E_zE_0\rangle_0),\qquad
 \Psi_n=n^{1+d}E_zS_0.
\]
Rounding the square radii leaves fixed buffers.  The two conditional
replacements preserve expectations in the plane and in every sufficiently
large plus disk: condition on both contours and use independence of their
interiors.  They consequently preserve charges by
Lemma~\ref{ref:charges}.  Proposition~\ref{ref:operator} gives
$\norm{\Phi_n}_\infty+\norm{\Psi_n}_\infty\le C$, with both supports in
$B_{Cn}(0)$.  Their scale-normalized charges are
\begin{equation}\label{ref:normalized-product-charges}
 \lambda_n=\frac{\ell(\Phi_n)}n=n\,\ell(e(z)e(0)),\qquad
 \mu_n=\frac{\ell_s(\Psi_n)}{n^d}=n\,\ell_s(e(z)\sigma_0),
\end{equation}
and are bounded uniformly in $n$.

For either $\Theta_n=\Phi_n$ or $\Psi_n$, let $b$ be its leading dimension,
and put $\chi_n=\ell_b(\Theta_n)/n^b$.  Subtract from $\Theta_n$ the field
$\chi_n n^bP_{Cn}\mathcal O_b$.  This subtraction has bounded norm and
zero charge.  Its expectation in a disk of radius $4Bn$, with $B$ fixed
and large, is therefore $O(B^{-p+\nu})$, by
\eqref{ref:improved-decay}.  Calibration of the subtracted line yields
\begin{equation}\label{ref:charge-extraction}
 \limsup_{n\to\infty}
 \left|\chi_n-\kappa_b^{-1}B^b d_{Bn}(\Theta_n)\right|
       \le CB^{b-p+\nu}.
\end{equation}
For a sequence of varying directions the same assertion follows from
compact-separated convergence after passing to a convergent subsequence
of $\zeta$.  This is the quantitative separation of the two limits: first
$n\to\infty$ with $B$ fixed, and then $B\to\infty$.

For the even product the limit of $d_{Bn}(\Phi_n)$ is
\[
 a^2\bigl(\langle\varepsilon(\zeta)\varepsilon(0)\rangle^+
                       _{\mathrm{disk}(4B)}-1\bigr).
\]
Rescale this disk by $B$.  The first expansion in
\eqref{ref:continuum-ope} gives
\[
 \langle\varepsilon(\zeta)\varepsilon(0)\rangle^+
                       _{\mathrm{disk}(4B)}
 =B^{-2}\langle\varepsilon(\zeta/B)\varepsilon(0)\rangle^+
                       _{\mathrm{disk}(4)}
 =1+O(B^{-2}).
\]
Thus $B d_{Bn}(\Phi_n)$ tends to zero in the iterated limit.  The error in
\eqref{ref:charge-extraction} also tends to zero, since the general even
exponent $p=2$ is already larger than $b=1$.  This proves
\eqref{ref:ee-charge}.

For the odd product, plane expectation vanishes, and the separated limit
is $aa_\sigma\langle\varepsilon(\zeta)\sigma(0)\rangle^+
_{\mathrm{disk}(4B)}$.  Let
$m_\sigma=\langle\sigma(0)\rangle^+_{\mathrm{disk}(4)}$.  The second
expansion in \eqref{ref:continuum-ope} gives
\[
 \langle\varepsilon(\zeta)\sigma(0)\rangle^+
                  _{\mathrm{disk}(4B)}
 =B^{-1-d}\bigl(\tfrac12 Bm_\sigma+O(1)\bigr).
\]
Since $\kappa_d=a_\sigma m_\sigma$, the odd instance of
\eqref{ref:charge-extraction} gives $\mu_n\to a/2$ in this sign
normalization.  Squaring gives \eqref{ref:es-charge}.  All the continuum
remainders just used are uniform for $|\zeta|=1$, which also proves the
claimed uniformity in direction.
\end{proof}

\subsection{Positive local tilts and finite boundary tests}

We next give two consequences needed for controlling sup norms of random
functions.  The first controls a large positive weight inside a small patch;
the second reduces a sup norm on a distant contour to finitely many scalar
evaluations.  In both, the untouched collar is in the reference model.

\begin{lemma}[Buffered comparison and normalized positive weights]
\label{ref:comparison}
Consider a finite critical reference graph in a square chart.  Let $I,J$ and
$V_{\mathrm{pin}}$ be pairwise disjoint vertex sets, and prescribe common
spin values on $V_{\mathrm{pin}}$.  Suppose $I\subset B_w(z)$ and
$J\subset B_R(z)^c$.  For $R\ge(1+c)w$ with fixed $c>0$, every assignment
$a$ on $I$ and assignments $\tau,\tau'$ on $J$ satisfy
\begin{equation}\label{ref:density-ratio}
 C_c^{-1}\le
 \frac{\Prob_0(\sigma_I=a\mid\sigma_J=\tau,\sigma_{V_{\mathrm{pin}}})}
      {\Prob_0(\sigma_I=a\mid\sigma_J=\tau',\sigma_{V_{\mathrm{pin}}})}
 \le C_c.
\end{equation}
For $R\ge8w\ge8$, the absolute logarithm of this ratio is at most
$C(w/R)^\theta$ for some $\theta>0$.  The constants are uniform in mesh,
common pins, and auxiliary finite real fields.  Either law can be replaced
by a mixture; the comparison also holds after a common conditional
randomization, across fixed polygonal collars of positive relative width,
and for different continuations whose reference interactions agree in the
separating chart.  Free or fixed physical boundaries can intersect such
collars.

In particular, let $W>0$ and $f$ be bounded functions supported in $B_w(z)$,
let $s/w>1$ leave a fixed relative collar, and put
\begin{equation}\label{ref:normalized-weights}
 Z=\langle W\rangle_0,\qquad
 N=Z^{-1}P_{s,z}W,\qquad M=Z^{-1}P_{s,z}(fW).
\end{equation}
Then
\begin{equation}\label{ref:weight-bounds}
 C^{-1}\le N\le C,\qquad \norm{M}_\infty\le C\norm{f}_\infty,
 \qquad \langle N\rangle_0=1,\quad
 \langle M\rangle_0=\frac{\langle fW\rangle_0}{Z}.
\end{equation}
The bounds are independent of the amplitude of $W$.  If $W$ and $f$ are
spin-flip even, then $N$ and $M$ are even.
\end{lemma}

\begin{proof}
The conditional-law assertions are the buffered and shrinking-patch
comparison \cite[Corollary~2.2]{ReferenceIsing}; the comparison across
physical boundaries has the same common-pin convention.  Apply
\eqref{ref:density-ratio} to the law in the inner square under an arbitrary
cutting configuration and to its mixture induced by the reference plane
law.  Pointwise comparison of these laws, integrated against $W\ge0$,
gives $C^{-1}Z\le P_{s,z}W\le CZ$.  The inequality
$|fW|\le\norm{f}_\infty W$ gives the bound for $M$.  The two mean
identities are the tower property.  Finally, global spin flip commutes
with the reference kernel, which proves the parity assertion.
\end{proof}

Taking $W=e^B$ in \eqref{ref:normalized-weights} makes both the denominator
and numerator of the tilted expectation bounded before a farther
reference cut is applied.  Lemma~\ref{norm:attenuation} will combine these
bounds with Proposition~\ref{ref:operator} to obtain attenuation uniform
in the size of an even local tilt $B$.

\begin{lemma}[Finite tests for an averaged sup norm]\label{ref:dictionary}
Fix concentric inner and outer squares with a sufficiently large fixed
ratio of radii and fixed relative buffers.  Rescale the inner radius to
one; the resulting mesh is $1/w$ when the original radius is $w$
in lattice units.  For every $\eta>0$ there are
an integer $N$ and a fine-mesh threshold such that, at each finer mesh,
one can choose outer cutting configurations $\tau_1,\ldots,\tau_N$ with
\begin{equation}\label{ref:finite-tests}
 \norm{PF}_\infty
      \le \max_{1\le k\le N}|PF(\tau_k)|+\eta\norm F_\infty
\end{equation}
for every bounded function $F$ supported in the inner square.  Here $P$
averages to the outer contour.  The number $N$ is independent of the mesh
and of $F$; the actual configurations may depend on the mesh.  Consequently,
for a random inner function $F$ and $q\ge1$,
\begin{equation}\label{ref:finite-moment-tests}
 \bigl(\E\norm{PF}_\infty^q\bigr)^{1/q}
 \le N^{1/q}\max_k\bigl(\E|PF(\tau_k)|^q\bigr)^{1/q}
       +\eta\bigl(\E\norm F_\infty^q\bigr)^{1/q}.
\end{equation}
\end{lemma}

\begin{proof}
Choose an intermediate annular contour separated from both squares by
fixed relative collars.  For any outer configuration $\tau$, the induced
law on this contour has a density $H_\tau$ relative to the reference
plane law satisfying $\langle H_\tau\rangle_0=1$ and
$0\le H_\tau\le C$.  This is the pointwise comparison used in
\cite[Proposition~4.1]{ReferenceIsing}.  Hence
$PF(\tau)=\langle FH_\tau\rangle_0$.

Apply the outer-to-inner version of finite soft-weight replacement
\cite[Lemma~3.3]{ReferenceIsing} to $H_\tau/C$.
For any $\eta_0>0$, it supplies finitely many positive normalized weights
$U_1,\ldots,U_m$ and vectors $c(\tau)\in\R^m$ with
$\sum_i|c_i(\tau)|\le C$, such that, simultaneously for all $F$,
\[
 \left|PF(\tau)-\sum_{i=1}^m c_i(\tau)
                                  \langle FU_i\rangle_0\right|
       \le C\eta_0\norm F_\infty.
\]
Because $U_i\ge0$ and $\langle U_i\rangle_0=1$, each functional
$F\mapsto\langle FU_i\rangle_0$ has norm at most one.  Cover the
finite-dimensional coefficient ball $\{c:\norm{c}_{\ell^1}\le C\}$ by finitely
many sets of $\ell^1$-diameter at most $\eta/2$.  From each occupied set
choose one actual outer configuration $\tau_k$.  For every $\tau$ in
that set, the two approximation errors and the coefficient difference give
\[
 |PF(\tau)-PF(\tau_k)|
       \le (2C\eta_0+\eta/2)\norm F_\infty.
\]
Take $\eta_0\le\eta/(4C)$.  The number of sets depends on the tolerance
and geometry, but not on the fine mesh.  Taking the supremum over $\tau$
proves \eqref{ref:finite-tests}.  Finally, Minkowski's inequality and
$\max_k|X_k|^q\le\sum_k|X_k|^q$ prove
\eqref{ref:finite-moment-tests}.
\end{proof}

The finite tests control images of arbitrary inner functions, so their
choice is independent of the disorder when $F$ is random.  Later we apply
them after summing all short inputs sent to one common parent contour.
That sum still has its support inside a fixed inner square; this is why
the buffer size in the blocking construction is chosen before the moment
estimates.
